\section{Objectifs du Projet}

Pour concevoir un système de détection de plagiat, nous avons défini plusieurs objectifs techniques et méthodologiques :

\subsection{Prétraitement et Extraction des Caractéristiques}
L’une des étapes clés du projet consiste à transformer un texte brut en une représentation exploitable par un algorithme de classification. Pour cela, nous avons :
\begin{itemize}
    \item[-] Développé un tokenizer capable de gérer la ponctuation, la casse, les stopwords et la lemmatisation.
    \item[-] Expérimenté différentes stratégies d’\textbf{n-grams} afin d’améliorer la capture des structures linguistiques.
    \item[-] Implémenté nos propres versions de \textbf{Bag-of-Words} et \textbf{TF-IDF}.
\end{itemize}

\subsection{Développement des Modèles de Classification}
Nous avons choisi d’implémenter plusieurs modèles afin de comparer leurs performances et d’identifier ceux les plus adaptés à notre problème. Nos implémentations incluent :
\begin{itemize}
    \item[-] \textbf{Naïve Bayes}, pour sa rapidité et son efficacité sur des données textuelles.
    \item[-] \textbf{Perceptron et Régression Logistique}, pour tester des approches linéaires d’apprentissage.
    \item[-] \textbf{Arbre de Décision, Random Forest et AdaBoost}, afin d’évaluer des approches non linéaires basées sur des ensembles de modèles.
\end{itemize}

\subsection{Optimisation et Expérimentations}
Une attention particulière a été portée à l’évaluation et à l’optimisation des modèles :
\begin{itemize}
    \item[-] Nous avons mis en place un \textbf{early stopping} pour éviter le sur-apprentissage et optimiser le temps d’entraînement.
    \item[-] Une évaluation détaillée des performances a été réalisée en utilisant les métriques \textbf{précision, rappel et F1-score}.
    \item[-] Nous avons analysé l'impact des hyperparamètres (taille des n-grams, régularisation, profondeur des arbres) sur les performances.
\end{itemize}

En suivant cette méthodologie, notre objectif est de construire une solution fiable et adaptable, tout en garantissant une compréhension approfondie des algorithmes sous-jacents.
